\section{\texorpdfstring{$k$-algebras, indecomposables, and idempotents}{k-algebras, indecomposables, and idempotents}}

\begin{definition}
    A nonzero $R$-module $M$ is called indecomposable if every decomposition
    \[
        M = M_1 \oplus M_2
    \]
    forces $M_1=0$ or $M_2=0$.
\end{definition}

\begin{theorem}[Krull-Schmidt] Let $M$ be an $R$-module with a composition series.
    Suppose
    \[
        M \cong V_1 \oplus \cdots \oplus V_m
        \cong
        U_1 \oplus \cdots \oplus U_n
    \]
    where all $V_i$ and $U_j$ are indecomposable.
    Then
    \[
        m=n,
    \]
    and after reordering the summands one has
    \[
        V_i \cong U_i
        \quad\text{for all } i.
    \]
\end{theorem}

\begin{proof}
    Replacing the given isomorphisms by identifications, we may assume
    \[
        M=V_1\oplus\cdots\oplus V_m
        =U_1\oplus\cdots\oplus U_n .
    \]
    If $M=0$, the statement is trivial. So assume $M\neq 0$.

    Since $M$ has a composition series, it has finite length. Hence each direct summand
    $V_i$ and $U_j$ also has finite length. Moreover, for every indecomposable module
    $X$ of finite length, Fitting's lemma implies that every endomorphism of $X$ is either
    invertible or nilpotent; therefore $\operatorname{End}_R(X)$ is a local ring.

    We prove the theorem by induction on the length $\ell(M)$ of $M$.
    If $\ell(M)=1$, then $M$ is simple, hence indecomposable, so necessarily $m=n=1$
    and $V_1\cong U_1$.

    Now assume $\ell(M)>1$ and that the statement is known for all modules of smaller
    length. Set
    \[
        A:=V_1,
        \qquad
        B:=V_2\oplus\cdots\oplus V_m,
    \]
    so that
    \[
        M=A\oplus B.
    \]
    Let
    \[
        \iota:A\to M,\qquad \pi:M\to A
    \]
    be the canonical inclusion and projection. For the decomposition
    \[
        M=U_1\oplus\cdots\oplus U_n,
    \]
    let
    \[
        \kappa_j:U_j\to M,\qquad p_j:M\to U_j
    \]
    be the canonical inclusion and projection for $1\le j\le n$.

    Since
    \[
        1_M=\sum_{j=1}^n \kappa_j p_j,
    \]
    we get
    \[
        1_A=\pi\iota
        =\pi\Bigl(\sum_{j=1}^n \kappa_j p_j\Bigr)\iota
        =\sum_{j=1}^n (\pi\kappa_j)(p_j\iota).
    \]
    For each $j$, put
    \[
        u_j:=(\pi\kappa_j)(p_j\iota)\in \operatorname{End}_R(A).
    \]
    If all $u_j$ were noninvertible, then, since $\operatorname{End}_R(A)$ is local, their sum would
    also be noninvertible. But their sum is $1_A$, a contradiction. Hence $u_j$ is a unit
    for some $j$. After reordering the $U_j$, we may assume that $u_1$ is a unit.

    Set
    \[
        a:=\pi\kappa_1:U_1\to A,
        \qquad
        b:=p_1\iota:A\to U_1.
    \]
    Then
    \[
        ab=u_1
    \]
    is an automorphism of $A$. Consider
    \[
        e:=b(ab)^{-1}a\in \operatorname{End}_R(U_1).
    \]
    Then
    \[
        e^2=b(ab)^{-1}a\,b(ab)^{-1}a
        =b(ab)^{-1}(ab)(ab)^{-1}a
        =b(ab)^{-1}a
        =e,
    \]
    so $e$ is an idempotent. Also,
    \[
        eb=b(ab)^{-1}ab=b\neq 0,
    \]
    hence $e\neq 0$. Since $\operatorname{End}_R(U_1)$ is local, its only idempotents are $0$ and $1$,
    therefore $e=1_{U_1}$.

    Thus
    \[
        b(ab)^{-1}a=1_{U_1},
    \]
    so $a$ is injective. On the other hand, $ab$ is surjective, hence $a$ is surjective as well.
    Therefore $a$ is an isomorphism, and we obtain
    \[
        A\cong U_1.
    \]

    Now define
    \[
        f:=\kappa_1 a^{-1}:A\to M.
    \]
    Then
    \[
        \pi f=\pi\kappa_1 a^{-1}=aa^{-1}=1_A.
    \]
    Hence
    \[
        M=f(A)\oplus B.
    \]
    Indeed, for any $x\in M$,
    \[
        x=f(\pi(x))+\bigl(x-f(\pi(x))\bigr),
    \]
    and
    \[
        \pi\bigl(x-f(\pi(x))\bigr)=\pi(x)-\pi f(\pi(x))=0,
    \]
    so $x-f(\pi(x))\in \ker\pi=B$. Also, if $f(a)\in B=\ker\pi$, then
    \[
        a=\pi f(a)=0,
    \]
    so the sum is direct. Since $f(A)=\kappa_1(U_1)=U_1$, we have
    \[
        M=U_1\oplus B.
    \]

    Therefore
    \[
        B\cong M/U_1\cong U_2\oplus\cdots\oplus U_n.
    \]
    But also
    \[
        B=V_2\oplus\cdots\oplus V_m.
    \]
    Now $\ell(B)<\ell(M)$, so by the induction hypothesis,
    \[
        m-1=n-1,
    \]
    and after reordering $U_2,\dots,U_n$, we have
    \[
        V_i\cong U_i \qquad (i=2,\dots,m).
    \]
    Together with $V_1\cong U_1$, this gives
    \[
        m=n,
        \qquad
        V_i\cong U_i \quad\text{for all } i
    \]
    after a suitable reordering.

    This completes the proof.
\end{proof}

\begin{definition}
    An idempotent of a ring $R$ is a nonzero element
    \[
        e \in R
    \]
    such that
    \[
        e^2=e.
    \]
\end{definition}

\begin{remark}
    If $e$ is an idempotent, then the corner ring
    \[
        eRe
    \]
    is a ring with identity element $e$.
\end{remark}

\begin{definition}
    Two idempotents $e_1,e_2 \in R$ are called orthogonal if
    \[
        e_1 e_2 = e_2 e_1 = 0.
    \]
    More generally, finitely many idempotents are orthogonal if they are pairwise orthogonal.
\end{definition}

\begin{lemma}
    If $e_1,\dots,e_n$ are orthogonal idempotents, then
    \[
        e_1+\cdots+e_n
    \]
    is again an idempotent.
\end{lemma}

\begin{proof}
    Let
    \[
        e:=e_1+\cdots+e_n.
    \]
    Since each $e_i$ is an idempotent, we have $e_i^2=e_i$ for all $i$. Since the
    $idempotents$ are orthogonal, for $i\neq j$ we have
    \[
        e_ie_j=0.
    \]
    Therefore
    \[
        e^2=(e_1+\cdots+e_n)^2
        =\sum_{i=1}^n e_i^2+\sum_{i\neq j} e_ie_j
        =\sum_{i=1}^n e_i+\sum_{i\neq j}0
        =e_1+\cdots+e_n
        =e.
    \]
    Hence $e$ is an idempotent.
\end{proof}

\begin{definition}
    An idempotent decomposition of an idempotent $e$ is an expression
    \[
        e=e_1+\cdots+e_n
    \]
    where the $e_i$ are orthogonal idempotents.
\end{definition}

\begin{remark}
    If $e=e_1+e_2$ is an idempotent decomposition, then $e_1$ and $e_2$ are in $eRe$.
\end{remark}

\begin{definition}
    An idempotent is called primitive if it cannot be written as a sum of two nonzero orthogonal idempotents.
\end{definition}

\begin{definition}
    A primitive orthogonal decomposition is an idempotent decomposition in which all summands are primitive.
\end{definition}

\begin{example}
    If $e$ is an idempotent in a unital ring, then
    \[
        1=e+(1-e)
    \]
    is an idempotent decomposition of $1$.
\end{example}

\begin{example}
    Let $e$ be an idempotent, $Re = \{ ae | a\in R\}$ is a left module of $R$.
\end{example}

\begin{theorem}
    Let $e$ be an idempotent of $R$.

    If
    \[
        e=e_1+\cdots+e_n
    \]
    is an idempotent decomposition, then
    \[
        Re=Re_1\oplus\cdots\oplus Re_n.
    \]

    Conversely, if $Re$ is a direct sum of left ideals
    \[
        Re=I_1\oplus\cdots\oplus I_n,
    \]
    then there exists an idempotent decomposition
    \[
        e=e_1+\cdots+e_n
    \]
    with $e_i\in I_i$ and $I_i=Re_i$.

    In this way, the idempotent decompositions of $e$ are in one-to-one correspondence
    with the direct sum decompositions of the left $R$-module $Re$.
\end{theorem}

\begin{proof}
    Assume first that
    \[
        e=e_1+\cdots+e_n
    \]
    is an idempotent decomposition. Since
    \[
        e_ie=ee_i=e_i,
    \]
    we have
    \[
        Re_i\subseteq Re.
    \]
    Indeed, if $x\in Re_i$, then $xe=x$.

    Now let $a\in R$. Then
    \[
        ae=ae_1+\cdots+ae_n\in Re_1+\cdots+Re_n,
    \]
    so
    \[
        Re\subseteq Re_1+\cdots+Re_n.
    \]
    The reverse inclusion is clear, hence
    \[
        Re=Re_1+\cdots+Re_n.
    \]

    To show that the sum is direct, suppose
    \[
        a_1e_1+\cdots+a_ne_n=0.
    \]
    Multiplying on the right by $e_i$, we get
    \[
        a_ie_i=0
    \]
    for each $i$, since $e_je_i=0$ for $j\ne i$. Therefore the expression is unique, and
    \[
        Re=Re_1\oplus\cdots\oplus Re_n.
    \]

    Conversely, suppose
    \[
        Re=I_1\oplus\cdots\oplus I_n.
    \]
    Since $e\in Re$, we may write
    \[
        e=e_1+\cdots+e_n,
        \qquad e_i\in I_i.
    \]

    For any $a_i\in I_i\subseteq Re$, we have
    \[
        a_i=a_ie=a_ie_1+\cdots+a_ie_n.
    \]
    Since $e_j\in I_j$, we have
    \[
        Re_j\subseteq I_j,
    \]
    hence
    \[
        a_ie_j\in I_j.
    \]
    But $a_i\in I_i$, and the decomposition
    \[
        Re=I_1\oplus\cdots\oplus I_n
    \]
    is direct. Therefore
    \[
        a_i=a_ie_i,
        \qquad
        a_ie_j=0 \quad (j\ne i).
    \]
    It follows that
    \[
        I_i\subseteq Re_i.
    \]
    On the other hand, since $e_i\in I_i$ and $I_i$ is a left ideal, we have
    \[
        Re_i\subseteq I_i.
    \]
    Thus
    \[
        I_i=Re_i.
    \]

    Now take $a_i=e_i$. Then
    \[
        e_i=e_i^2,
        \qquad
        e_ie_j=0 \quad (j\ne i).
    \]
    Hence
    \[
        e=e_1+\cdots+e_n
    \]
    is an idempotent decomposition.
\end{proof}

\begin{corollary}
    Let $e$ be an idempotent of $R$. Then the following are equivalent:
    \begin{enumerate}
        \renewcommand{\labelenumi}{(\theenumi)}
        \item $e$ is primitive.
        \item $Re$ is indecomposable as a left $R$-module.
        \item The only idempotents in $eRe$ are $0$ and $e$.
    \end{enumerate}
\end{corollary}

\begin{proof}
    \textbf{(1) \(\Longleftrightarrow\) (2).}
    This follows immediately from the previous theorem, which gives a bijection between
    idempotent decompositions of $e$ and direct sum decompositions of the left ideal $Re$.

    \medskip

    \textbf{(1) \(\Longrightarrow\) (3).}
    Suppose that $e' \in eRe$ is an idempotent with
    \[
        e' \neq 0, \qquad e' \neq e.
    \]
    Since $e' \in eRe$, we have
    \[
        ee' = e' = e'e.
    \]
    Hence
    \[
        e = e' + (e-e').
    \]
    We claim that this is an idempotent decomposition of $e$.

    First,
    \[
        (e-e')^2 = e^2 - ee' - e'e + (e')^2
        = e - e' - e' + e'
        = e - e'.
    \]
    So $e-e'$ is an idempotent.

    Next,
    \[
        e'(e-e') = e'e - (e')^2 = e' - e' = 0,
    \]
    and similarly
    \[
        (e-e')e' = ee' - (e')^2 = e' - e' = 0.
    \]
    Thus $e'$ and $e-e'$ are orthogonal idempotents. Since both are nonzero, this contradicts
    the primitiveness of $e$.

    Therefore no such $e'$ exists, and the only idempotents in $eRe$ are $0$ and $e$.

    \medskip

    \textbf{(3) \(\Longrightarrow\) (1).}
    Suppose that $e$ is not primitive. Then there exist nonzero orthogonal idempotents
    $e_1,e_2$ such that
    \[
        e = e_1 + e_2.
    \]
    Since
    \[
        ee_1e = e_1, \qquad ee_2e = e_2,
    \]
    we have
    \[
        e_1,e_2 \in eRe.
    \]
    But $e_1$ is a nonzero idempotent different from $e$, contradicting (3). Hence $e$ must be
    primitive.
\end{proof}

